\begingroup
\footnotesize
\SingleSpacing
\setlength{\tabcolsep}{2pt}
\renewcommand{\UrlLeft}{}
\renewcommand{\UrlRight}{}
\newsavebox{\caixaseparadorvariaveis}
\sbox{\caixaseparadorvariaveis}{%
    \vbox to 1.25pt{\vss\color{black!25}\hrule width \linewidth height 0.35pt\vss}}
\newcounter{linhavariavelAdoze}
\newcommand{\marcalinhavariavel}{%
    \stepcounter{linhavariavelAdoze}%
    \label{tab:a12-linha-\arabic{linhavariavelAdoze}}}
% O espaço é constante para que as referências de página convirjam.
\newcommand{\separadorvariaveis}{\marcalinhavariavel\\\noalign{%
    \ifnum\getpagerefnumber{tab:a12-linha-\arabic{linhavariavelAdoze}}=%
          \getpagerefnumber{tab:a12-linha-\number\numexpr\value{linhavariavelAdoze}+1\relax}\relax
        \cleaders\copy\caixaseparadorvariaveis
    \fi
    \vskip 1.25pt\relax}}
\begin{longtable}{@{}>{\raggedright\arraybackslash}p{0.23\linewidth}
                     >{\raggedright\arraybackslash}p{0.18\linewidth}
                     >{\raggedright\arraybackslash}p{0.29\linewidth}
                     >{\raggedright\arraybackslash}p{0.12\linewidth}
                     >{\raggedright\arraybackslash}p{\dimexpr0.18\linewidth-8\tabcolsep\relax}@{}}
\caption{Dicionário das variáveis dos modelos da Seção~\ref{sec:discussao}}
\label{tab:ap-dicionario-h1}\\
\toprule
\textbf{Nome no código} & \textbf{Variável} & \textbf{Descrição} & \textbf{Unidade de medida} & \textbf{Comentário} \\
\midrule
\endfirsthead
\caption[]{Dicionário das variáveis (continuação)}\\
\toprule
\textbf{Nome no código} & \textbf{Variável} & \textbf{Descrição} & \textbf{Unidade de medida} & \textbf{Comentário} \\
\midrule
\endhead
\midrule
\endfoot
\bottomrule
\endlastfoot

\multicolumn{5}{@{}l}{\textit{Variáveis dependentes e controle dinâmico}} \\*
\nolinkurl{y_agr} & Conversão para agricultura \newline $Y_{it}^{\text{agr}}$
    & Área que passou de vegetação nativa em $t-1$ para agricultura em $t$, dividida pelo estoque nativo em $t-1$ e multiplicada por 100.
    & p.p. do estoque nativo inicial
    & MapBiomas 10.1. Numerador e denominador usam a mesma máscara de observação.\separadorvariaveis
\nolinkurl{y_pas} & Conversão para pastagem \newline $Y_{it}^{\text{past}}$
    & Área que passou de vegetação nativa em $t-1$ para pastagem em $t$, dividida pelo estoque nativo em $t-1$ e multiplicada por 100.
    & p.p. do estoque nativo inicial
    & MapBiomas 10.1. A taxa é de conversão bruta, assim como no destino agrícola.\separadorvariaveis
\nolinkurl{milho_plant_pct} & Participação da área plantada de milho \newline $A_{it}^{\text{milho}}$
    & Área municipal plantada com milho no ano $t$, dividida pela área do município e multiplicada por 100.
    & \% da área municipal
    & Área plantada da PAM/IBGE.\separadorvariaveis
\nolinkurl{y_w} & Conversão agrícola winsorizada
    & Taxa $Y_{it}^{\text{agr}}$ limitada superiormente ao percentil 99 da amostra observada.
    & p.p. do estoque nativo inicial
    & Valores acima do limite são substituídos pelo percentil.\separadorvariaveis
\nolinkurl{y_agr} em \nolinkurl{df_v2} & Conversão agrícola na grade completa
    & Taxa $Y_{it}^{\text{agr}}$ na grade de 3.384 pares município--ano, incluindo valores preenchidos na base precária.
    & p.p. do estoque nativo inicial
    & Dos 612 pares sem conversão observada, 594 têm preenchimento positivo.\separadorvariaveis
\nolinkurl{y_agr_l1}, \nolinkurl{y_l1} & Conversão agrícola defasada \newline $Y_{i,t-1}^{\text{agr}}$
    & Taxa de conversão para agricultura observada no mesmo município no ano anterior.
    & p.p. do estoque nativo inicial
    & Controle de persistência no modelo dinâmico de H1b.\separadorvariaveis
\multicolumn{5}{@{}l}{\textit{Preços internacionais e tendência}} \\*
\nolinkurl{lps} & Preço real da soja \newline $\ln P_{s,t-1}$
    & Logaritmo natural centrado da média anual do preço internacional real da soja em $t-1$.
    & log (adim.)
    & Série em US\$ de 2024/t antes do logaritmo.\separadorvariaveis
\nolinkurl{lpm} & Preço real do milho \newline $\ln P_{m,t-1}$
    & Logaritmo natural centrado da média anual do preço internacional real do milho em $t-1$.
    & log (adim.)
    & Série em US\$ de 2024/t antes do logaritmo.\separadorvariaveis
\nolinkurl{lpb} & Preço real da carne bovina \newline $\ln P_{\text{boi},t-1}$
    & Logaritmo natural centrado da média anual do preço internacional real da carne bovina em $t-1$.
    & log (adim.)
    & Série em US\$ de 2024/kg antes do logaritmo.\separadorvariaveis
\nolinkurl{lcb_s} & Preço relativo soja/boi \newline $\ln(P_s/P_{\text{boi}})_{t-1}$
    & Diferença entre os logaritmos centrados dos preços reais da soja e da carne bovina, ambos em $t-1$.
    & log (adim.)
    & Termo \textit{crop-to-beef} nas equações de agricultura e pastagem.\separadorvariaveis
\nolinkurl{log(ps_lead)-c_ps}\newline\nolinkurl{log(pm_lead)-c_pm}
    & Preço futuro da soja ou do milho \newline $\ln P_{j,t+1}$
    & Logaritmo natural do preço real médio no ano seguinte, centrado pela mesma média do preço defasado.
    & log (adim.)
    & Placebo temporal.\separadorvariaveis
\nolinkurl{lps_reg}, \nolinkurl{lpm_reg} & Preço centrado por regime
    & Logaritmo do preço real defasado da soja ou do milho, menos a média de seu logaritmo no regime correspondente.
    & log (adim.)
    & Painel B da Tabela~\ref{tab:h1-regimes-biomas}.\separadorvariaveis
\nolinkurl{trend} & Tendência temporal $\tau_t$
    & Tempo decorrido desde 2001, calculado como $(t-2001)/10$.
    & década
    &\separadorvariaveis
\multicolumn{5}{@{}l}{\textit{Aptidão e ocupação prévia}} \\*
\nolinkurl{zs}, \nolinkurl{zm} & Aptidão de alto insumo \newline $Z_{i,j}$
    & Rendimento potencial municipal de soja ou milho no cenário de alto insumo e sequeiro do GAEZ v4.
    & escore $z$
    & Rendimento original em kg/ha. Padronização nos 141 municípios.\separadorvariaveis
\nolinkurl{gaez_yxLr0_soy_soja} & Aptidão de baixo insumo
    & Rendimento potencial municipal da soja no cenário de baixo insumo e sequeiro do GAEZ v4.
    & escore $z$
    & Rendimento original em kg/ha.\separadorvariaveis
\nolinkurl{gaez_sxHr0_soy_soja} & Índice de adequação da soja
    & Índice de adequação do GAEZ v4 para soja em alto insumo e sequeiro.
    & escore $z$
    & Cenário de índice \texttt{sxHr0}.\separadorvariaveis
\nolinkurl{pam_1991_2000_min7_soja} & Rendimento histórico da soja
    & Rendimento médio municipal da soja registrado na PAM de 1991 a 2000.
    & escore $z$
    & Rendimento original em kg/ha. Exige ao menos sete anos válidos.\separadorvariaveis
\nolinkurl{z_exp_agr} & Histórico agrícola $H_i$
    & Média das frações anuais do município ocupadas por lavouras temporárias entre 1985 e 2000.
    & escore $z$
    & MapBiomas 10.1.\separadorvariaveis
\nolinkurl{z_exp_soja} & Histórico de soja $H_i$
    & Fração da área municipal ocupada por soja em 2000.
    & escore $z$
    & Medida alternativa de ocupação prévia na Tabela~\ref{tab:h1-margens-tecnologia}.\separadorvariaveis
\nolinkurl{z_exp_pas} & Histórico pastoril \newline $H_i^{\text{past}}$
    & Média das frações anuais do município ocupadas por pastagem entre 1985 e 2000.
    & escore $z$
    & MapBiomas 10.1.\separadorvariaveis
\multicolumn{5}{@{}l}{\textit{Clima, logística e fiscalização}} \\*
\nolinkurl{rain} & Anomalia de precipitação \newline $R_{it}$
    & Precipitação anual municipal em $t$ menos a média anual de 1981--2000.
    & 100 mm
    & CHIRPS. Valor positivo indica um ano mais chuvoso que a média histórica.\separadorvariaveis
\nolinkurl{prox_sul_z} & Proximidade portuária \newline $P_i^{\text{sul}}$
    & Negativo do escore $z$ da menor distância geodésica do centroide municipal a Santos ou Paranaguá.
    & escore $z$
    & Valores maiores indicam maior proximidade.\separadorvariaveis
\nolinkurl{dist_min_t_100km} & Distância portuária mínima \newline $D_{it}^{\min}$
    & Menor distância geodésica a Santos ou Paranaguá até 2013. A partir de 2014, o cálculo também considera Miritituba.
    & 100 km
    & Medida indireta do custo de acesso aos mercados.\separadorvariaveis
\nolinkurl{choc_ferrov_13} & Exposição à ferrovia \newline $F_{it}^{2013}$
    & Proximidade ao terminal ferroviário de Rondonópolis, multiplicada pelo indicador de 2013 em diante.
    & escore $z$
    & Proximidade é o negativo do escore $z$ da distância geodésica.\separadorvariaveis
\nolinkurl{f1}\newline\nolinkurl{inst} (de \nolinkurl{autos_l1})
    & Autos de Flora \newline $F_{i,t-1}$
    & $\ln(1+n)$, em que $n$ é o total de autos de infração de Flora do IBAMA no município em $t-1$.
    & log (adim.)
    & Usa registros geolocalizados e com classificação temática completa na subamostra fiscal.\separadorvariaveis
\nolinkurl{emb_area_l1}\newline\nolinkurl{inst} (de \nolinkurl{area_l1})
    & Área de novos embargos \newline $E_{i,t-1}^{\text{ha}}$
    & $\ln(1+a)$, em que $a$ é a soma das áreas, em hectares, dos novos termos de embargo do IBAMA no município em $t-1$.
    & log (adim.)
    & Termos geolocalizados, deduplicados e sem cancelamento registrado.\separadorvariaveis
\nolinkurl{emb_esp_l1}\newline\nolinkurl{inst} (de \nolinkurl{termos_l1})
    & Contagem de novos embargos \newline $E_{i,t-1}^{\text{n}}$
    & $\ln(1+n)$, em que $n$ é o número de novos termos de embargo do IBAMA no município em $t-1$.
    & log (adim.)
    & Mesmos critérios de seleção da medida por área.\separadorvariaveis
\nolinkurl{emb_area_f1}, \nolinkurl{emb_esp_f1} & Embargos futuros \newline $E_{i,t+1}^{\text{ha}}$, $E_{i,t+1}^{\text{n}}$
    & Medidas de área e contagem de novos embargos construídas da mesma forma, usando o ano seguinte.
    & log (adim.)
    & Placebo de fiscalização.\separadorvariaveis
\multicolumn{5}{@{}l}{\textit{Território, regimes e efeitos fixos}} \\*
\nolinkurl{bioma} & Bioma predominante $b(i)$
    & Bioma com a maior fração territorial no município (Amazônia, Cerrado ou Pantanal).
    & categoria
    & Define os recortes e os efeitos fixos de bioma por ano.\separadorvariaveis
\nolinkurl{amz}, \nolinkurl{pan} & Frações territoriais de bioma \newline $q_i^M$, $q_i^{\text{Pan}}$
    & Frações da área municipal situadas na Amazônia e no Pantanal, respectivamente.
    & fração (0 a 1)
    & O Cerrado é a referência no teste de moderação contínua.\separadorvariaveis
\nolinkurl{d_p1}, \nolinkurl{d_p2}, \nolinkurl{d_p3} & Indicadores de regime $I_{rt}$
    & Indicadores dos períodos 2001--2005, 2006--2011 e 2012--2024. Cada indicador assume 1 no respectivo período.
    & binária (0/1)
    &\separadorvariaveis
\nolinkurl{post05} & Indicador do degrau de 2005 \newline $D_{2005}$
    & Indicador que assume 1 a partir de 2005, inclusive.
    & binária (0/1)
    & Coluna (3) da Tabela~\ref{tab:h1-destinos-robustez}.\separadorvariaveis
\nolinkurl{is_agr} & Indicador do destino agrícola \newline $I_{\text{agr}}$
    & Indicador que assume 1 para agricultura e 0 para pastagem no painel empilhado.
    & binária (0/1)
    & Cada par município--ano entra com os dois destinos.\separadorvariaveis
\nolinkurl{ibge}, \nolinkurl{ano}, \nolinkurl{bioma^ano}\newline
\nolinkurl{ibge^dest}, \nolinkurl{bioma^ano^dest}
    & Efeitos fixos de município, de ano e de bioma por ano
    & Categorias de município (código IBGE), ano civil ou par bioma predominante--ano.
    & categoria
    & No modelo empilhado, os efeitos de município e de bioma por ano variam por destino.\separadorvariaveis
\multicolumn{5}{@{}l}{\textit{Interações e termos derivados}} \\*
\nolinkurl{xs}, \nolinkurl{xm}\newline
\nolinkurl{xs_yxLr}, \nolinkurl{xs_sxHr}, \nolinkurl{xs_pam90}\newline
\nolinkurl{xs_f}, \nolinkurl{xm_f}
    & Preço $\times$ aptidão
    & Produto de $\ln P_{j,t-1}$ por $Z_{i,j}$, ou por uma medida alternativa de aptidão.
    & log $\times$ escore $z$
    & No placebo, usa-se o preço em $t+1$.\separadorvariaveis
\nolinkurl{x_exp_agr_s}, \nolinkurl{x_exp_agr_m}, \nolinkurl{x_exp_soja_s}\newline
\nolinkurl{x_pb_zpas}
    & Preço $\times$ ocupação prévia
    & Produto do preço dos grãos com $H_i$, ou do preço da carne bovina com $H_i^{\text{past}}$.
    & log $\times$ escore $z$
    & Os preços entram em log e defasados.\separadorvariaveis
\nolinkurl{zs_tr}, \nolinkurl{zm_tr}\newline
\nolinkurl{zs_yxLr_tr}, \nolinkurl{zs_sxHr_tr}, \nolinkurl{zs_pam90_tr}\newline
\nolinkurl{z_exp_agr_tr}, \nolinkurl{z_exp_soja_tr}, \nolinkurl{zpas_tr}
    & Tendência $\times$ aptidão ou ocupação prévia
    & Produto de $\tau_t$ pela medida de aptidão ou de ocupação prévia usada na especificação.
    & década $\times$ escore $z$
    & Permite tendências distintas por aptidão ou ocupação.\separadorvariaveis
\nolinkurl{x_pb_zs}, \nolinkurl{x_pb_zm}, \nolinkurl{x_cb_s_zs}
    & Preço da carne bovina ou preço relativo $\times$ aptidão
    & Produto de $\ln P_{\text{boi},t-1}$ ou $\ln(P_s/P_{\text{boi}})_{t-1}$ pelo escore GAEZ da cultura.
    & log $\times$ escore $z$
    &\separadorvariaveis
\nolinkurl{x_ps_proxsul}, \nolinkurl{x_pm_proxsul}\newline
\nolinkurl{xs_proxsul}, \nolinkurl{xm_proxsul}
    & Preço e aptidão $\times$ proximidade portuária
    & Produtos de $P_i^{\text{sul}}$ pelo log do preço defasado e pelo gradiente preço--aptidão.
    & log $\times$ um ou dois escores $z$
    & Inclui a interação dupla e a tripla.\separadorvariaveis
\nolinkurl{f1_ps}, \nolinkurl{f1_pm}, \nolinkurl{f1_zs}, \nolinkurl{f1_zm}, \nolinkurl{f1_xs}, \nolinkurl{f1_xm}\newline
\nolinkurl{emb_ps}, \nolinkurl{emb_pm}, \nolinkurl{emb_zs}, \nolinkurl{emb_zm}, \nolinkurl{emb_xs}, \nolinkurl{emb_xm}\newline
\nolinkurl{emb_area_ps}, \nolinkurl{emb_area_pm}, \nolinkurl{emb_area_zs}, \nolinkurl{emb_area_zm}, \nolinkurl{emb_area_xs}, \nolinkurl{emb_area_xm}\newline
\nolinkurl{emb_ps_f1}, \nolinkurl{emb_zs_f1}, \nolinkurl{emb_xs_f1}\newline
\nolinkurl{inst_preco}, \nolinkurl{inst_aptidao}, \nolinkurl{inst_exposicao}
    & Preço e aptidão $\times$ fiscalização
    & Interações do log do preço defasado e do escore de aptidão com autos de Flora, área embargada ou contagem de embargos.
    & log $\times$ log, log $\times$ escore $z$ ou produto dos três fatores
    & Inclui os termos duplos e triplos usados em cada modelo e a interação com a contagem futura no placebo.\separadorvariaveis
\nolinkurl{xs_amz}, \nolinkurl{xs_pan}, \nolinkurl{xm_amz}, \nolinkurl{xm_pan}
    & Preço--aptidão $\times$ fração de bioma
    & Produto do gradiente preço--aptidão pelas frações municipais de Amazônia e Pantanal.
    & log $\times$ escore $z$ $\times$ fração
    & Permite que a inclinação varie com a composição territorial.\separadorvariaveis
\nolinkurl{xs_p1}, \nolinkurl{xs_p2}, \nolinkurl{xs_p3}\newline
\nolinkurl{xm_p1}, \nolinkurl{xm_p2}, \nolinkurl{xm_p3}\newline
\nolinkurl{xs_r_p1}, \nolinkurl{xs_r_p2}, \nolinkurl{xs_r_p3}\newline
\nolinkurl{xm_r_p1}, \nolinkurl{xm_r_p2}, \nolinkurl{xm_r_p3}\newline
\nolinkurl{x_reg2}, \nolinkurl{x_reg3}
    & Preço--aptidão $\times$ regime
    & Produto do gradiente preço--aptidão por cada indicador $I_{rt}$.
    & log $\times$ escore $z$
    & Usa preço centrado dentro do regime no Painel B da Tabela~\ref{tab:h1-regimes-biomas}.\separadorvariaveis
\nolinkurl{zs_p2}, \nolinkurl{zs_p3}, \nolinkurl{zm_p2}, \nolinkurl{zm_p3}
    & Aptidão $\times$ regime
    & Produto de $Z_{i,j}$ pelos indicadores de 2006--2011 e 2012--2024.
    & escore $z$
    & Deslocamentos de nível em relação a 2001--2005.\separadorvariaveis
\nolinkurl{e_amz}, \nolinkurl{e_z}, \nolinkurl{e_zamz} & Degrau de 2005
    & Produtos de $D_{2005}$ pela fração amazônica, pela aptidão e pelo produto dessas duas medidas.
    & fração, escore $z$ ou produto
    & A tabela de robustez reporta o termo com aptidão.\separadorvariaveis
\nolinkurl{xs_agr}, \nolinkurl{xm_agr}, \nolinkurl{zs_tr_agr}, \nolinkurl{zm_tr_agr}, \nolinkurl{rain_agr}
    & Interações por destino
    & Produtos de $I_{\text{agr}}$ com o gradiente preço--aptidão, a tendência por aptidão e a anomalia de precipitação.
    & unidade do fator interagido
    & Permitem comparar a resposta na agricultura com a resposta na pastagem.\marcalinhavariavel \\
\end{longtable}
\noindent\textbf{Nota.} A primeira coluna reproduz os nomes ou as expressões usados no notebook de estimação e no script complementar da revisão da qualificação. Os índices $i$ e $t$ indicam município e ano. O índice $j$ identifica a commodity e $s$ e $m$ designam soja e milho nas interações com aptidão. As séries de preços são da \textit{Pink Sheet} do Banco Mundial, deflacionadas pelo CPI-U para dólares de 2024. Na especificação de referência, os logaritmos são centrados pela média da série defasada no painel de 2001--2024. As unidades referem-se à escala usada na regressão. A abreviação p.p. indica pontos percentuais e adim. indica uma escala adimensional.
\par
\endgroup
